\chapter{Expanding curves that smooth a corner}\label{sec:examples-expanders}
Expanding solitons give nonstationary examples defined for positive times and converging to a nonsmooth initial cone. We construct an even convex graphical family directly. The primary comparison is \cite{Hal12}, Section 6, where expanding planar curves are studied in arc-length coordinates. Here the normalization is $\vec H=x^\perp/2$, so the time family is $\sqrt t\,\Sigma$.

\begin{proposition}[The expander graph equation and its positive defect]\label{prop:expander-graph-defect}
A graph $y=u(x)$ in $\mathbb R^2$ satisfies the scale-one expander equation exactly when
\[
 u''=\frac12(1+(u')^2)(u-xu').
\]
If $u(0)=a>0$ and $u'(0)=0$, set $p=u'$ and $q=u-xp$. On every interval of existence with $x\geq0$,
\[
 p'=\tfrac12(1+p^2)q,\qquad q'=-xp',\qquad
 q(x)=a\exp\left(-\frac12\int_0^x s(1+p(s)^2)\,ds\right)>0.
\]
In particular $p'>0$ and $0<q(x)\leq ae^{-x^2/4}$.
\end{proposition}
\begin{proof}
For the upward normal $\nu=(-p,1)/\sqrt{1+p^2}$, the coefficient of the mean curvature vector along $\nu$ is $u''/(1+p^2)^{3/2}$, while $\langle(x,u),\nu\rangle=(u-xp)/\sqrt{1+p^2}$. Their equality with factor one half gives the ODE. Differentiate $q=u-xp$ to obtain $q'=-xp'=-x(1+p^2)q/2$. Solving this scalar linear equation with $q(0)=a$ gives its positive exponential expression and Gaussian upper bound. The scalar convention $H=\operatorname{div}\nu$ is the negative of the coefficient of $\vec H$ here.
\end{proof}

\begin{lemma}[The expanding graph has bounded slope]\label{lem:expander-global-slope}
For each $a>0$, the initial-value solution in Proposition \ref{prop:expander-graph-defect} extends smoothly to all $x\geq0$. For every $\varepsilon>0$ in its interval of existence and $x\geq\varepsilon$,
\[
 0<p(x)\leq p(\varepsilon)+\frac{q(\varepsilon)}\varepsilon.
\]
It therefore has a finite positive limiting slope $\ell_a=\lim_{x\to\infty}p(x)$.
\end{lemma}
\begin{proof}
The smooth ODE has a unique local solution for these initial data. Since $p'>0$ and $q>0$, integration of $q'=-xp'$ gives
\[
 q(\varepsilon)-q(x)=\int_\varepsilon^x s p'(s)\,ds
 \geq\varepsilon(p(x)-p(\varepsilon)).
\]
This proves the bound. Were the maximal endpoint finite and positive, choose a fixed positive $\varepsilon$ below it. The bound controls $p$, while $0<q\leq a$ controls $u=q+xp$ on that finite interval. Smooth ODE continuation then crosses the endpoint, a contradiction. Monotonicity and boundedness of $p$ give its finite limit, which is positive since $p(x)>0$ for $x>0$.
\end{proof}

\begin{theorem}[A complete convex expander asymptotic to two rays]\label{thm:convex-graphical-expander}
For every $a>0$, there is a unique smooth even solution $u_a:\mathbb R\to\mathbb R$ of the expander graph equation with $u_a(0)=a$ and $u_a'(0)=0$. Its graph is properly embedded, complete, and strictly convex. For a finite $\ell_a>0$,
\[
 u_a(x)-\ell_a|x|\longrightarrow0\quad(|x|\to\infty),
 \qquad 0<u_a(x)-\ell_a|x|\leq a.
\]
Its two asymptotic rays meet at the origin. This constructs a family; it does not assert that a prescribed cone angle determines a unique expander among arbitrary immersions.
\end{theorem}
\begin{proof}
The ODE is invariant under $x\mapsto-x$, so local uniqueness makes the solution even; Lemma \ref{lem:expander-global-slope} extends both halves. Positivity of $q$ gives $u_a''>0$. As an entire smooth graph it is proper, embedded, and complete. Write $\ell=\ell_a$. Boundedness of $p$ and the Gaussian bound for $q$ give $0<p'(x)\leq Ce^{-x^2/4}$, hence for $x\geq1$,
\[
 0<\ell-p(x)\leq C\int_x^\infty e^{-s^2/4}\,ds
 \leq\frac{2C}{x}e^{-x^2/4}.
\]
Since $u_a-x\ell=q-x(\ell-p)$, both terms tend to zero. Its derivative is $p-\ell<0$ for $x\geq0$; thus it decreases from $a$ to zero and is strictly positive. Reflecting gives the negative half. These properties agree with the expanding-curve models in \cite{Hal12}, Section 6, after converting its dilation normalization.
\end{proof}

\begin{proposition}[A smooth flow converging uniformly to a corner]\label{prop:expander-corner-initial-trace}
Let $u_a$ be the expander above and define
\[
 U(x,t)=\sqrt t\,u_a(x/\sqrt t),\qquad t>0.
\]
Its graphs solve curve shortening flow, and for every $x\in\mathbb R$,
\[
 0<U(x,t)-\ell_a|x|\leq a\sqrt t.
\]
Thus the graph functions converge uniformly on all of $\mathbb R$ to the Lipschitz corner $\ell_a|x|$ as $t\downarrow0$. Their products with $\mathbb R^k$ give hypersurface flows with an initial crease.
\end{proposition}
\begin{proof}
Put $\xi=x/\sqrt t$. Direct differentiation gives $U_t=(u_a-\xi u_a')/(2\sqrt t)$, $U_x=u_a'$, and $U_{xx}=u_a''/\sqrt t$. The expander equation makes $U_t=U_{xx}/(1+U_x^2)$. The bound follows by multiplying the height excess in Theorem \ref{thm:convex-graphical-expander} by $\sqrt t$. Proposition \ref{prop:example-soliton-product} handles products. This is an explicit classical solution for $t>0$ with an initial trace; no uniqueness theorem for arbitrary weak flows from that trace is used.
\end{proof}

\begin{proposition}[Total curvature and maximum curvature of the expander]\label{prop:expander-curvature-example}
The expander graph has finite total curvature
\[
 \int_\Sigma|\kappa|\,ds=2\arctan\ell_a<\pi.
\]
Its maximum curvature is $a/2$, attained at its tip. Along its expanding flow the maximum is $a/(2\sqrt t)$.
\end{proposition}
\begin{proof}
Since $u_a''>0$, the curvature magnitude is $p'/(1+p^2)^{3/2}$ and $ds=\sqrt{1+p^2}\,dx$. Their product is $(\arctan p)'dx$. The limiting slopes are $-\ell_a$ and $\ell_a$, proving the integral. By the equation, $|\kappa|=q/(2\sqrt{1+p^2})$. On $x>0$, $q$ strictly decreases and $p$ strictly increases, so the curvature decreases from $q(0)/2=a/2$. Evenness handles the other half. Dilation by $\sqrt t$ divides curvature by $\sqrt t$.
\end{proof}
